\documentclass{article}
\usepackage[utf8]{inputenc}
\begin{document}
\section{Exercise 4: Keep It Short}

This exercise is intended to teach strategies for identifying problems with flow, transition and organization across writing.

Your tasks for this exercise are the following:
\begin{itemize}
\item Read the following paragraphs. Remove any unnecessary phrases, vague comments and make it as concise as possible. 
\item There are a total of 209 words. How low can you go?
\item (Side note: this is an example of ineffective writing. While no grammatical errors are being made, the author constantly switches between first and third person, uses colloquialism and generally rambles on. This is not a model for scholarly writing!)
\end{itemize}

Note: there is no need for additions or intense editing (grammar/spelling) - at least for this exercise! If you are running into any problems with Overleaf or have questions, please feel free to raise your hand and ask for help.

\subsection*{Sample Paragraph}

As we’ve all understood by now, circadian rhythms are approximately twenty-four hour rhythms that are driven by intracellular clocks. These consist of transcriptional/translational feedback loops. However, intercellular communication is also important for circadian rhythms – which is what my essay today will discuss. 

Before I get to the nitty gritty details, I’d like to introduce the concept of the circadian clock, in terms of mammals. The circadian clock can be found in the suprachiasmatic nucleus (SCN). We'll be focusing largely on \textit{Per2}, so let’s walk through what feedback loops \textit{Per2} is involved in.

An example of a positive feedback loop with \textit{Per2} is when the BMAL1/CLOCK heterodimer directly activates the transcription of the \textit{Per2} gene. Additionally, the circadian clock controls downstream events by regulating the expression of CCGs – clock controlled genes. We’ll come across these later in the paper.

Circadian rhythms may have a role in growth control – it hasn’t been solidly established yet in previous literature. This leads us to the study’s objective.

We want to investigate the role of circadian genes in growth control and DNA damage response \textit{in vivo}. To do this, we will use mice. We specifically used an \textit{mPer2} mutant. These mice are homozygous for the mPer2 mutation and are deficient in circadian clock function. 


ewpage
\subsection*{Suggested Techniques}

There are multiple ways to approach this exercise. A few suggestions are given below. 
\paragraph{} Note: this is not a complete list.

\begin{itemize}
\item There is no need for "As we've all understood by now" or "which is what my essay today will discuss."
\item Twenty-four can be rewritten as 24. As a general rule, numbers between 1-10 are written as words. Anything larger than 10 should be written as numerals.
\item This can be removed: "Before I get to the nitty gritty details, I’d like to introduce the concept of..". There's no point in announcing what you're about to do to the reader - \textit{do it instead}.
\item Can be more concise by stating that "In mammals, the circadian clock can be found in the..."
\item The purpose can be rewritten to exclude the use of \textit{we}.
\item The second last paragraph can be easily summarized with one sentence. "Previous literature suggests that circadian rhythms may play a role in growth control."
\item The last three statements can be combined. E.g. "In this study, the \textit{mPer2} mutant mice were used, as the homozygous \textit{mPer2} mutation results in them being deficient in circadian clock function."
\end{itemize}


\end{document}